% Generated draft: requires mathematical and source review.
\documentclass[11pt]{article}
\usepackage[margin=25mm]{geometry}
\usepackage{fontspec}
\setmainfont{Latin Modern Roman}
\usepackage{amsmath,amssymb,mathtools,unicode-math}
\setmathfont{Latin Modern Math}
\usepackage{enumitem,needspace,xurl,hyperref}
\setlength{\parindent}{0pt}
\setlength{\parskip}{5pt}
\setlength{\emergencystretch}{3em}
\newcommand{\N}{\mathbb{N}}
\newcommand{\Z}{\mathbb{Z}}
\newcommand{\Q}{\mathbb{Q}}
\newcommand{\R}{\mathbb{R}}
\newcommand{\C}{\mathbb{C}}
\newcommand{\F}{\mathbb{F}}
\newcommand{\A}{\mathbb{A}}
\newcommand{\PP}{\mathbb{P}}
\newcommand{\E}{\mathbb{E}}
\newcommand{\eps}{\varepsilon}
\newcommand{\dd}{\,\mathrm{d}}
\begin{document}
\section*{UnsolvedMath ID: 11000280}
\textbf{Source problem code:} catalogue entry
\subsection{Definitions and mathematical statement}
Let \(F_3\) be the free group of rank \(3\), and let
\[
 \operatorname{Out}(F_3)=\operatorname{Aut}(F_3)/\operatorname{Inn}(F_3).
\]
A complex linear representation is a group homomorphism into \(\operatorname{GL}_m(\C)\), and is faithful when its kernel is trivial.

\emph{Question 2.12.} Does there exist a finite positive integer \(m\) and an injective homomorphism
\[
 \rho:\operatorname{Out}(F_3)\hookrightarrow\operatorname{GL}_m(\C)?
\]
The rank is exactly \(3\), and the coefficient field is exactly \(\C\). The source imposes no dimension bound, discreteness of the image, or additional geometric condition.

The source notes braid-group linearity [8] and says closed-surface mapping-class-group linearity was unknown in its context. It credits Brendle and Hamidi--Tehrani [13] with showing that the Formanek--Procesi approach cannot be directly adapted there: the particular ``poison subgroup'' described in earlier, unprovided text does not arise. This does not prove that all obstructions to linearity are absent or that those mapping class groups are linear.

It describes the persistence of the \(\operatorname{Out}(F_3)\) question as evidence of behavior different from larger ranks and mentions the existence of finite-index subgroups mapping onto \(\Z\) as another such instance. The exact comparison ranks and hypotheses are not supplied. Its promised automatic-structure and isoperimetric discussion is absent. All these are source-era comparisons, not new claims of current resolution.
\subsection{Short English statement}
Is the rank-three outer automorphism group of a free group linear over \(\C\) in some finite dimension? The source's braid, mapping-class, and larger-rank comparisons are retained as context without changing the fixed-rank question.
\subsection{Sources}
\begin{itemize}
\item Local UnsolvedMath registry, permanent ID 11000280, source code catalogue entry, Question 2.12, PDF page 332. No individual proposer is supplied.
\item Farb, \emph{Problems on Mapping Class Groups and Related Topics}. Supplied URL: \url{https://www.math.uchicago.edu/~farb/papers/mcgbook.pdf}; recorded access 2026-07-29; extraction pdf-font-structure. The volume contains 329 explicitly typeset headings across chapters, with repeated numbering disambiguated by page.
\item Reference [8] is cited for braid linearity; Brendle and Hamidi--Tehrani [13] for the poison-subgroup obstruction comparison; Formanek and Procesi are named without a full citation. The earlier poison definition is not supplied.
\item Supplied status open and OPEN-TRIAGE. The research summary names Grunewald--Lubotzky, gives explicitly uncertain statements about larger-rank linearity, and says it cannot verify the status. Those uncertain remarks are not adopted as results. Current primary-literature status and redistribution terms remain NEEDS\_REVIEW.
\end{itemize}
\end{document}
